\section{Discussion and limitations}
\label{sec:discussion}
\paragraph{What the evidence says.}
With inputs fixed, the room to shorten PPI++ is set by how far $E[Y\mid f,X]$ is from affine in $f$, which can be estimated before choosing a method. On the seven real pools it is at most about 1.2\% for the mean for the adjuster classes tried (a lower bound on the true gap), and all but nine of 270 adjuster cells are wider than PPI++ (Sec.~\ref{sec:real}),
so a practitioner with these data should keep PPI++ and, if intervals undercover at small $n$, calibrate them (Sec.~\ref{sec:validity}) rather than add an adjuster. The cases where an adjuster does pay off share two features: a large gap $\eta^2-R^2_{\rm aff}$, and $n$ beyond a threshold set by the adjuster's size ($n=50$ already for a $(d{+}2)$-parameter linear map, a few hundred for boosting).
A cheap diagnostic suggested by the population calculation of Sec.~\ref{sec:headroom}, which we did not evaluate as a decision rule: fit the candidate adjusters by cross-fitting on the labelled set and compare their out-of-fold $R^2$ with that of the affine fit; if the gain in $R^2/(1+\rho)$ is below the estimation cost, expect a loss.
Gains also fade as $\rho$ grows (Sec.~\ref{sec:theory}): a large unlabelled set makes PPI++ harder to beat (Sec.~\ref{sec:real}).

\paragraph{Limitations.}
\emph{Sampling protocol.} Real-pool experiments resample the pool i.i.d.\ with replacement and take the pool mean (or fit) as truth; this is a clean finite population and not a fresh draw from a super-population, and $N\le2\times10^4$ throughout, so small $\rho$ is covered only by the synthetic study and by the sweep at $n=200$ (Sec.~\ref{sec:real}).
\emph{Assumptions.} Labels are missing completely at random, labelled and unlabelled records are exchangeable; under the $2013\to2014$ shift on \texttt{plankton} (Appendix~\ref{app:diag}) PPI++ and the two adjusters tested lose coverage and nothing in this paper repairs it.
\emph{Scope.} Five adjuster classes with default hyper-parameters, $K=5$ folds, mean and OLS only (no quantiles or logistic regression), a binary predictor only in the diagnostics; a better-tuned or richer learner could shift the break-even $n$ but not the floor.
\emph{Theory.} The efficiency floor is a restated, known bound; the estimation-cost argument is first-order and we do not prove a break-even $n$ nor test it quantitatively. \emph{Synthetic scenarios} have headroom by construction and show what is possible, not what is typical.
\emph{Validity.} Wald, $t$ and the bootstrap undercover at small $n$ on skewed or rare-outcome pools (down to 0.755 at $n=15$), and ratios are only read where coverage is comparable; we claim no finite-sample guarantee for any of the intervals.
\emph{Statistical power.} Ratios carry paired confidence intervals; a ratio within 0.5\% of 1 is called parity or a negligible effect, not a finding.
Intervals that are shorter with no loss of coverage are a benefit; we see no harm beyond the stated misuse of narrower-but-undercovering intervals. The census pools are public, anonymised microdata; age and sex enter only as regression covariates, and we infer or release nothing about individuals.
